\chapter{Market-Making Games}\label{ch:m2:market-making-games}

An interviewer points out of the window: make me a market on the number of windows in
that building. You say 200 at 400. She says: I buy. You have just sold her windows at 400,
and you have learned something: she thinks there are more than 400, or at least she is
willing to bet so. Your next market should be higher. Trading firms train and test
people with games like this because every market maker's problem is in it: estimate a
value you cannot see, quote around it wide enough to survive the people who know more,
narrow enough to trade with the people who do not, and learn from every trade. This
chapter plays the game seriously: how to make a market, how to update on trades, how
adverse selection sets the width, and what firms say about the games they use.

\section{Make me a market}

\begin{definition}[Make-me-a-market, width]\label{def:m2:market-making-games:mmam}
\emph{Make-me-a-market}\index{make-me-a-market} is a request to quote, on some uncertain
quantity, a bid at which one will buy and an offer at which one will sell, for a stated
size, before knowing which way the other party will trade. The \emph{width}\index{width}
of the market is the offer minus the bid.
\end{definition}

The language is that of an open-outcry pit, and Jane Street's public guide to its trading
interviews spells it out: ``I'm 2 bid for 10'' offers to buy 10 at 2; ``I have 10 at 4''
offers to sell; ``2 at 4, 10 up'' does both; the other side says ``sold'' to hit the bid or
``take 'em'' to lift the offer, and the order stands until it trades or its maker says
``I'm out''. The guide's advice for a contract paying the roll of a die, worth 3.5 on
average, is to buy below and sell above the expected value, to weigh the worst loss
against one's capital, and to balance the chance of trading against the profit per
trade; it suggests ``3 at 4, 10 up''.

A market has a centre, the maker's estimate of the value, and a width. The centre is a
statistics problem; the width is an economic one. Quote too narrow and anyone who knows
the value better trades against you when you are wrong; quote too wide and no one
trades. The card game of this chapter makes both precise.

\begin{definition}[Inventory]\label{def:m2:market-making-games:inv}
A market maker's \emph{inventory}\index{inventory} is the net position it holds from the
trades done against its quotes, long after net buying and short after net selling.
\end{definition}

\section{Updating on trades}

\begin{definition}[Bayesian update]\label{def:m2:market-making-games:bayes}
A \emph{Bayesian update}\index{Bayesian update} revises a probability distribution over an
unknown value after an observation, in proportion to how likely the observation would be
under each value: posterior $\propto$ prior $\times$ likelihood.
\end{definition}

The contract of the chapter's game pays the sum of five cards dealt face down from a
52-card deck, aces 1 to kings 13. Before any trade the sum is worth 35 on average, with a
standard deviation of 8.03, and it lies between 6 and 64
(\cref{fig:m2:market-making-games:dist}). Each round a trader arrives: with probability
0.2 an informed one, who knows the sum and buys if it is above the offer, sells if below
the bid, and otherwise passes; otherwise an uninformed one, who trades on a random side
with a probability that falls with the width, $e^{-w/8}$, and otherwise passes.

\begin{omfigure}
\begin{tikzpicture}
\begin{axis}[width=11cm,height=4.6cm,ybar,bar width=2.2pt,
  xlabel={sum of five cards},ylabel={probability},
  tick label style={font=\small},label style={font=\small},
  xmin=4,xmax=66,ymin=0,ymax=0.055,grid=major,grid style={gray!25},scaled y ticks=false,
  yticklabel style={/pgf/number format/fixed,/pgf/number format/precision=2},
  table/col sep=comma]
\addplot[fill=omDef!60,draw=omDef] table[x=sum,y=p]{figdata/markets-2/30-market-making-games/distribution.csv};
\addplot[omThm,very thick,dashed,sharp plot] coordinates {(35,0) (35,0.055)};
\end{axis}
\end{tikzpicture}

\omcaption{Distribution of the sum of five cards dealt from a 52-card deck, aces counting 1
and kings 13: mean 35 (dashed), standard deviation 8.03, from 6 to 64. Exact, by
counting. Data: the chapter's tutorial.}\label{fig:m2:market-making-games:dist}
\end{omfigure}

\begin{example}[The first trade]\label{ex:m2:market-making-games:first}
The maker opens 30 at 40, a width of 10. An uninformed trader trades with probability
$e^{-10/8} = 0.29$. If someone buys at 40, it was an informed trader with probability 0.31;
the maker's expected value of the sum rises from 35 to 38.13, and its next market is
centred there. After a sale at 30 it falls to 31.87; after a pass it stays at 35, since a
pass is equally likely from a high or a low sum.
\end{example}

Over a game the maker's centre drifts towards the true sum as trades reveal it
(\cref{fig:m2:market-making-games:game}); how fast depends on the share of informed
traders and on the width, since a wide market filters out the informed trades that carry
the information along with the uninformed ones that pay for the width.

\begin{omfigure}
\begin{tikzpicture}
\begin{axis}[width=11cm,height=4.6cm,
  xlabel={round},ylabel={maker's mid},
  tick label style={font=\small},label style={font=\small},
  xmin=0,xmax=20,ymin=28,ymax=42,grid=major,grid style={gray!25},
  legend columns=2,legend style={font=\footnotesize,at={(0.5,-0.3)},anchor=north,draw=none,fill=none,/tikz/every even column/.append style={column sep=8pt}},
  table/col sep=comma]
\addplot[omDef,very thick,const plot,mark=*,mark size=1.2pt] table[x=round,y=mid]{figdata/markets-2/30-market-making-games/game.csv};
\addplot[omThm,thick,dashed] table[x=round,y=value]{figdata/markets-2/30-market-making-games/game.csv};
\legend{maker's expected value,true sum (31)}
\end{axis}
\end{tikzpicture}

\omcaption{One game at a width of 11.5: the maker's expected value of the sum before each
round, updated by Bayes' rule after every trade and pass, against the true sum, 31. Early
buys push the estimate up; the sales that follow bring it down towards the truth. The
maker ends the game with a profit of 18.6. Data: the chapter's tutorial,
seeded.}\label{fig:m2:market-making-games:game}
\end{omfigure}

\section{Adverse selection around a table}

The guide's definition of adverse selection matches the one in One Quant Book 1,
chapter 1: the trades you get to do are worse than they would appear, because someone else is
selecting into the trade against you. Its practical form is a question to ask before
every quote: conditional on this order being filled, what is the thing worth, and would I
still want the trade? Glosten and Milgrom turned it into a model in 1985: a specialist facing
informed and uninformed traders sets its bid and offer at the expected value conditional
on a sale and on a purchase, and the presence of better-informed traders produces a
positive spread even when the specialist is risk-neutral and makes zero expected profit.

\begin{definition}[Winner's curse]\label{def:m2:market-making-games:curse}
The \emph{winner's curse}\index{winner's curse} is the tendency of the winner of an auction
among bidders with noisy estimates of a common value to have overestimated it: winning is
itself evidence that one's estimate was too high.
\end{definition}

A market maker is always the winner of such an auction: its quote was the one someone
chose to trade on. The width is the price of the curse, and the best width balances two
effects (\cref{fig:m2:market-making-games:pnl}). A wider market loses less to the
informed, who trade only when the sum lies outside it, but trades less with the
uninformed, who pay half the width each time.

\begin{proposition}[Width against the table]\label{prop:m2:market-making-games:width}
With no informed traders, the maker's expected profit per round is $\tfrac w2 e^{-w/\lambda}$
with patience $\lambda$, maximised at $w = \lambda$. Informed traders make every width
cost more and the narrowest widths cost most, since they trade only against quotes on the
wrong side of the value; learning from trades lowers the loss to them, and so allows a
narrower market.
\end{proposition}

\begin{proof}
An uninformed trader trades with probability $e^{-w/\lambda}$ and, the centre being the
expected value, earns the maker $w/2$ on average; the product is maximised at $w = \lambda$.
An informed trader's expected cost to the maker is $E[(V - a)^+ + (b - V)^+]$, which falls
as the quotes move apart; a better centre shrinks it for any width.
\end{proof}

\begin{example}[The best width]\label{ex:m2:market-making-games:best}
Over 2\,000 seeded games of 20 rounds with a fifth of the traders informed, the maker who
updates earns most, 16.4 per game, at a width of 11.5 points, against 14.1 at 8 points and
15.1 at 15. A maker who never updates does best wider, at 14 points, and earns only 12.0.
Without informed traders the best width would be the patience, 8 points, earning about
29 a game (29.4 by the formula, 28.6 in the simulation).
\end{example}

\begin{omfigure}
\begin{tikzpicture}
\begin{axis}[width=11cm,height=5cm,
  xlabel={width of the market (points)},ylabel={maker's P\&L per game},
  tick label style={font=\small},label style={font=\small},
  xmin=0,xmax=30,ymin=-30,ymax=25,grid=major,grid style={gray!25},
  legend columns=4,legend style={font=\footnotesize,at={(0.5,-0.3)},anchor=north,draw=none,fill=none,/tikz/every even column/.append style={column sep=6pt}},
  table/col sep=comma]
\addplot[omDef,very thick] table[x=width,y=learn10]{figdata/markets-2/30-market-making-games/pnl.csv};
\addplot[omThm,very thick] table[x=width,y=learn20]{figdata/markets-2/30-market-making-games/pnl.csv};
\addplot[omProp,very thick] table[x=width,y=learn30]{figdata/markets-2/30-market-making-games/pnl.csv};
\addplot[omThm,thick,dashed] table[x=width,y=naive20]{figdata/markets-2/30-market-making-games/pnl.csv};
\addplot[gray,thin,forget plot] coordinates {(0,0) (30,0)};
\legend{10\% informed,20\%,30\%,20\%: no updating}
\end{axis}
\end{tikzpicture}

\omcaption{The maker's expected P\&L per game of 20 rounds against the width of its
market, for tables with 10, 20 and 30\% informed traders, updating by Bayes' rule, and for
20\% without updating (dashed). Narrow markets are picked off; wide ones do not trade; the
best width rises, and the profit falls, with the share of informed traders. 600 seeded
games per point. Data: the chapter's tutorial.}\label{fig:m2:market-making-games:pnl}
\end{omfigure}

\section{The training games firms describe}

Firms rarely publish their games, but some describe them. Jane Street's guide sets
practice markets: on the roll of a twenty-sided die, the maximum of three six-sided dice,
a die with the option to reroll once, the temperature in an hour, and a million times a
die roll, with the question of how one would feel if one's bid were hit and the roll came
up one. Its trading internship lists mock trading, in teams, on scenarios built by its
traders, and classes that range from financial products to poker. The common thread is the
chapter's: the value is uncertain, the other players know different things, and every
trade is information.

The games also teach inventory. A maker who has sold three times in a row is short, and
if its sales were to informed traders the value is above where it sold. Shifting both
quotes up, or skewing them, protects it, at the cost of trading less on one side; the
card engine can be extended to do so.

\section{Tutorial: the card-sum market}\label{tut:m2:market-making-games:mm}

\begin{tutorial}
\textbf{Goal.} Compute the exact distribution of the card sum, play the market against a
table of informed and uninformed bots with and without Bayesian updating, and find the
width that maximises expected profit.
\textbf{End state:} \cref{fig:m2:market-making-games:dist,fig:m2:market-making-games:game,fig:m2:market-making-games:pnl},
\cref{ex:m2:market-making-games:first,ex:m2:market-making-games:best} and the numbers of
the weekend problem.
\begin{tutsteps}
\item \textbf{The deck and the update}: the exact distribution of the sum, the table's
parameters, and Bayes' rule after a trade or a pass.
\omcode{code/firm/mmgame/firm_mmgame.py}{15}{49}{The card-sum distribution, the table, and the Bayesian update.}\label{lst:m2:market-making-games:update}
\item \textbf{The game}: dealing and playing one game.
\omcode{code/firm/mmgame/firm_mmgame.py}{52}{80}{Dealing and one game.}\label{lst:m2:market-making-games:play}
\item \textbf{The widths}: the expected P\&L of each width against the same deals.
\omcode{code/firm/mmgame/firm_mmgame.py}{83}{95}{Expected P\&L by width, over common deals.}\label{lst:m2:market-making-games:pnl}
\item \textbf{Run} \texttt{mmgame\_demo.first\_trade()}, \texttt{mmgame\_demo.problem()}
and \texttt{fig\_mmgame.py}.
\end{tutsteps}
\textbf{What to change next.} Let the maker skew its quotes with its inventory, and find
whether it earns more; then give each bot one of the five cards instead of the sum, so
that every trader is partly informed, as around a real table.
\end{tutorial}

\section{Build: the market-making game engine}\label{bld:m2:market-making-games:mmgame}

\begin{build}
\bfield{Purpose} The miniature firm trains its quoting logic, and its people, on games
whose answers it can compute.
\bfield{Interface} \texttt{card\_sum\_distribution(k, copies, values)};
\texttt{Table(cards, rounds, informed, patience)}; \texttt{mean(post)};
\texttt{update(post, action, bid, ask, table, width)}; \texttt{deal(table, rng, prior)};
\texttt{play(width, table, value, draws, learn, prior)}; \texttt{expected\_pnl(widths,
table, games, seed, learn)}.
\bfield{Rules} Five cards from one deck; one trader a round; informed traders trade only
when the sum is outside the quotes; uninformed ones trade with probability
$e^{-w/\lambda}$ on a random side; quotes centred on the maker's current expected value;
every width faces the same seeded deals.
\bfield{Acceptance tests} \texttt{code/firm/mmgame/tests/}: the distribution sums to one
with mean 35 and the right extremes; a buy raises and a sale lowers the estimate; without
informed traders the profit matches the formula; updating beats not updating; games are
reproducible.
\bfield{Stretch} Inventory skew; several makers competing on width; traders with partial
information (one card each); the Glosten--Milgrom zero-profit spread; a live version for
people to play against the bots.
\end{build}

\begin{omsources}
\item Jane Street, ``Probability and markets guide''; trading internship page.
\item L.~R.~Glosten and P.~R.~Milgrom, ``Bid, ask and transaction prices in a specialist
market with heterogeneously informed traders'', Journal of Financial Economics 14 (1985).
\end{omsources}

\section{Exercises}

\begin{exercise}[$\star$]\label{exo:m2:market-making-games:1}
Make a market on the roll of a six-sided die, and explain your width.
\end{exercise}

\begin{exercise}[$\star$]\label{exo:m2:market-making-games:2}
You quote 200 at 400 on the number of windows and the interviewer buys. What do you do
with your next market, and why?
\end{exercise}

\begin{exercise}[$\star$]\label{exo:m2:market-making-games:3}
Why does a pass leave the maker's estimate at 35 on the first round of the card game?
\end{exercise}

\begin{exercise}[$\star\star$]\label{exo:m2:market-making-games:4}
With no informed traders and patience 8, what width maximises the maker's profit per
round, and what does it earn over 20 rounds?
\end{exercise}

\begin{exercise}[$\star\star$]\label{exo:m2:market-making-games:5}
In the first round at 30 at 40, what is the probability that a buyer is informed? Why is
it so much higher than the 20\% of informed traders at the table?
\end{exercise}

\begin{exercise}[$\star\star$]\label{exo:m2:market-making-games:6}
Explain the \omterm{def:m2:market-making-games:curse}{winner's curse} in a market maker's terms.
\end{exercise}

\begin{exercise}[$\star\star\star$]\label{exo:m2:market-making-games:7}
\emph{Coding.} With \texttt{expected\_pnl}, find the best width for a table with 30\%
informed traders on the grid 8 to 15 in half points, over the same 2\,000 seeded deals.
\end{exercise}

\begin{exercise}[$\star\star\star$]\label{exo:m2:market-making-games:8}
\emph{Find the flaw.} ``I quote tight so I trade more: more trades means more spread
earned.'' Correct it.
\end{exercise}

\section{Problem: The Card-Sum Market}

\begin{problem}[{Weekend problem --- the width that pays}]\label{pb:m2:market-making-games:1}
You make markets for 20 rounds on the sum of five cards dealt face down from a 52-card
deck. Each round one trader comes: a fifth of the time one who knows the sum; otherwise
one who trades on a random side with probability $e^{-w/8}$ at a width $w$. You quote
around your current expected value.

\medskip\noindent\textbf{Part I --- The contract.}
\begin{enumerate}
\item What are the mean, standard deviation and range of the sum?
\item Why are the extremes 6 and 64 rather than 5 and 65?
\item At a width of 10, how likely is an uninformed trader to trade?
\item Quote your first market at width 10.
\item What do you earn per round, on average, from an uninformed trader?
\end{enumerate}

\medskip\noindent\textbf{Part II --- Learning.}
\begin{enumerate}[resume]
\item After a buy at 40, what is your new expected value? After a sale at 30?
\item Why does a pass teach you nothing on the first round, but something later?
\item In the chapter's example game, how does your estimate move, and why does it first
go the wrong way?
\item What does updating earn you, at the best widths, against not updating?
\item Why does the maker who does not update prefer a wider market?
\end{enumerate}

\medskip\noindent\textbf{Part III --- The width.}
\begin{enumerate}[resume]
\item Find the best width and the expected P\&L per game.
\item What do widths of 8 and 15 earn?
\item Without informed traders, what would the best width and profit be?
\item How does the best width change with 10\% and 30\% informed traders?
\item How likely does an uninformed trader trade at the best width?
\end{enumerate}

\medskip\noindent\textbf{Part IV --- Judgement.}
\begin{enumerate}[resume]
\item How would you know, at a real table, what share of traders is informed?
\item How should you adjust your quotes after three sales in a row?
\item What does this game teach about quoting real securities?
\item State the \emph{named result}: the width that maximises expected P\&L against this
table, and that P\&L.
\item In one sentence: what does the width pay for?
\end{enumerate}
\end{problem}

\section{Interview questions}

\begin{interviewq}[$\star$ \iqroles{trader}]\label{iq:m2:market-making-games:1}
Make me a market on the sum of two dice.
\end{interviewq}

\begin{interviewq}[$\star$ \iqroles{trader}]\label{iq:m2:market-making-games:2}
I buy on your offer. What do you do now?
\end{interviewq}

\begin{interviewq}[$\star\star$ \iqroles{trader, researcher}]\label{iq:m2:market-making-games:3}
Why does a bid--ask spread exist even with no costs? Explain with a model.
\end{interviewq}

\begin{interviewq}[$\star\star$ \iqroles{trader}]\label{iq:m2:market-making-games:4}
You are short three contracts after three sales. How do you move your market?
\end{interviewq}

\begin{interviewq}[$\star\star$ \iqroles{researcher}]\label{iq:m2:market-making-games:5}
How would you estimate the share of informed traders in a market from its trades?
\end{interviewq}

\begin{interviewq}[$\star\star\star$ \iqroles{developer}]\label{iq:m2:market-making-games:6}
Design a multi-player market-making game server for training, with bots and humans.
\end{interviewq}
